% SPDX-License-Identifier: Apache-2.0
% covers: PerBridge
\datasheet{PerBridge}{The peripheral side of an exit: request packets in, AXI to a peripheral out.}

\dsfacts{\code{lib/parts/src/bus/noc/bridge.rs}}{\code{txhdl\_parts::bus::noc::bridge::PerBridge}}%
{\code{//docs:noc}}

\subsection*{Function}

\code{PerBridge} is the other end of \code{HostBridge}. It takes
request packets on \code{req} from its node's exit and unpacks them
into \code{aw}, \code{ar} and \code{w} for a peripheral's tracker. It
takes the tracker's answers on \code{b} and \code{r}, and packs them on
\code{rsp} back to the node each request came from.

It also renames. An identifier belongs to the host that made it, so
two hosts may use the same one. The bridge gives every request a local
identifier of its own, and remembers against it the source node and
the identifier that node used. The answer comes back under the local
identifier, and the bridge sends the packet home under the original.

\subsection*{Parameters}

\begin{center}\footnotesize
\begin{tabular}{@{}lp{0.7\textwidth}@{}}
\toprule
Parameter & Meaning \\
\midrule
\code{X}, \code{Y} & The column and row of the bridge's node. The
bridge stamps them on every answer as its source. \\
\code{XB}, \code{YB} & The widths of a column and a row coordinate, in
bits. \\
\code{A} & The address width of the AXI link. \\
\code{D} & The data width of the AXI link. \\
\code{S} & The strobe width, which is \code{D / 8}. \\
\code{I} & The identifier width of the AXI link. \\
\code{NIDS} & How many local identifiers the bridge gives out, and so
how many bursts the peripheral has in flight at most. \\
\bottomrule
\end{tabular}
\end{center}

The tables below are for \code{PerBridge<0, 1, 2, 2, 32, 32, 4, 2,
4>}, the bridge of the memory at 0,1 in \code{soc/src/lib.rs}.

\dstables{PerBridge}

\subsection*{Behaviour}

\begin{itemize}
\item The local identifier on offer is \code{turn}. It is free when its
bit in \code{busy} is low.
\item A request packet with \code{chan} \code{W} is a write. The bridge
takes it when \code{turn} is free and both \code{aw} and \code{w} are
ready. It sends the address phase on \code{aw} under the local
identifier, and the data and strobe on \code{w} with \code{last} high.
\item The bridge treats any other request packet as a read. It takes it
when \code{turn} is free and \code{ar} is ready, and sends the address
phase on \code{ar} under the local identifier.
\item On taking a request, the bridge stores the source column in
\code{sx}, the source row in \code{sy} and the original identifier in
\code{oid}, each at the local identifier. It sets that identifier's
\code{busy} bit and adds one to \code{turn}.
\item A request waits while the identifier whose turn it is has not
been answered. The bridge does not look for another free identifier.
\item The bridge sends at most one answer packet per cycle, when
\code{rsp} has room, and a write response before a read beat. The
packet goes to the \code{sx} and \code{sy} stored under the answer's
identifier, under the \code{oid} stored there, with the source
\code{X}, \code{Y}.
\item A write response goes with \code{chan} \code{B}, its response,
and \code{last} high. A read beat goes with \code{chan} \code{R}, its
data, its response and its \code{last}.
\item The bridge clears an identifier's \code{busy} bit when it sends a
write response, or the last beat of a read, under that identifier.
\end{itemize}

\subsection*{Verification}

\begin{itemize}
\item Three tests in \code{lib/parts/src/bus/noc.rs}, in
\code{//lib/parts:parts\_test}, put a
\code{PerBridge<1, 1, 2, 2, 16, 32, 4, 2, 4>} in front of the memory
at 1,1.
\code{a\_burst\_crosses\_the\_lattice\_and\_its\_answer\_comes\_back}
writes and reads one word from one host.
\code{two\_hosts\_share\_a\_memory\_and\_each\_answer\_goes\_home} checks
that each answer reaches the host that asked. \code{//docs:noc} states
that this test failed before the renaming was written.
\code{four\_peripherals\_share\_one\_node} puts a router of four
memories behind the bridge, and checks each word and the router's
\code{DecErr} for an address none of them has.
\item \code{a\_core\_draws\_a\_solid\_and\_a\_gpu\_fills\_it}, in
\code{//soc:demo\_test}, runs a bridge at 0,1 for the memory and one at
1,1 for the serial port.
\item No waveform in \code{docs/BUILD.bazel} lowers the bridge, so its
netlist has no co-simulation under nvc or Verilator.
\code{//docs:showcase} lists the network's netlists as not yet
simulated.
\end{itemize}

\subsection*{Used by}

\code{soc::run} in \code{soc/src/lib.rs}: at 0,1 in front of the
memory's \code{AxiPer}, and at 1,1 in front of the serial port, behind
\code{LiteBridge1}. The tests of the network in
\code{lib/parts/src/bus/noc.rs}.

\subsection*{Limits}

A write request is one packet with one beat, so a burst of more than
one beat does not cross the network. At most \code{NIDS} bursts are in
flight at the peripheral. \code{turn} is a register of \code{I} bits
and \code{busy} has \code{NIDS} bits. Every use in the tree sets
\code{NIDS} to 4 with \code{I} of 2. The bridge sends one answer per
cycle.
